\subsection{Betrag von Summen und Differenzen}
\label{subsec:vector_algebra_magnitude_sums_differences}

Das Skalarprodukt macht sichtbar, wie der eingeschlossene Winkel den Betrag
von $\vec a+\vec b$ beziehungsweise $\vec a-\vec b$ beeinflusst.

\begin{center}
  \CachedTikzFigure{content/resources/chapters/vector_algebra/tikz/vector_sum_difference_magnitudes.tex}
\end{center}

Für $\vec a,\vec b\ne\vec 0$ und den eingeschlossenen Winkel $\varphi$ gilt:
\begin{ClozeArea}[34]{7}
  \[
    \begin{aligned}
      |\vec a\pm\vec b|^2
      &=(\vec a\pm\vec b)\cdot(\vec a\pm\vec b)\\[0.5em]
      &=\vec a\cdot\vec a+\vec b\cdot\vec b
        \pm2\,\vec a\cdot\vec b\\[0.5em]
      &=|\vec a|^2+|\vec b|^2
        \pm2|\vec a|\,|\vec b|\cos(\varphi).
    \end{aligned}
  \]
\end{ClozeArea}

\begin{KeyPointBox}{Beträge von Vektorkombinationen}{vector_algebra_sum_difference_magnitudes}
  \[
    \boxed{
      |\vec a\pm\vec b|
      =\sqrt{|\vec a|^2+|\vec b|^2
        \pm2|\vec a|\,|\vec b|\cos(\varphi)}
    }.
  \]
\end{KeyPointBox}

\begin{KeyPointBox}{Kosinussatz als Differenzfall}{vector_algebra_cosine_rule_derivation}
  Bilden $\vec a$ und $\vec b$ zwei Dreiecksseiten mit gemeinsamem
  Anfangspunkt und ist $\vec c=\vec a-\vec b$, so folgt unmittelbar
  \[
    \boxed{c^2=a^2+b^2-2ab\cos(\varphi)}.
  \]
  Für $\varphi=90^\circ$ bleibt der Satz des Pythagoras.
\end{KeyPointBox}

\begin{ExampleBox}{Differenz zweier gleich langer Vektoren}{vector_algebra_difference_magnitude_example}
  Für $|\vec a|=|\vec b|=\qty{3}{\centi\meter}$ und
  $\varphi=30^\circ$ gilt
  \[
    |\vec a-\vec b|^2
    =(\qty{3}{\centi\meter})^2+(\qty{3}{\centi\meter})^2
    -2(\qty{3}{\centi\meter})(\qty{3}{\centi\meter})\cos(30^\circ),
  \]
  \[
    |\vec a-\vec b|
    =\Cloze{5cm}{\sqrt{18-9\sqrt3}\,\unit{\centi\meter}
      \approx\qty{1.55}{\centi\meter}}.
  \]
\end{ExampleBox}
